% !TEX program = pdflatex
% ------------------------------------------------------------
% MATH 521 Homework Template (adapted from your MATH 421 HW10)
%
% HOW TO USE (quick):
%   1) Edit the Metadata block (HW number / due date especially).
%   2) Paste each problem statement under \problem[]{<n>}.
%   3) Write in the Scratch/Discussion box if you want.
%   4) Write your proof under "Proof." and end with \qed (or \qedhere).
%
% Notes:
% - Keeps theorem/definition boxes (tcolorbox).
% - Keeps ToC + PDF bookmarks for each problem.
% - Adds a Scratch/Discussion box per problem (optional).
% - Header bar has course + meeting info (left) and HW number (right).
% ------------------------------------------------------------

% TROUBLESHOOTING NOTE:
% If you ever see errors like `\tcb@insert@after@part` or other tcolorbox internals being
% an "Undefined control sequence", that usually means LaTeX is accidentally loading a
% DIFFERENT (often older) `tcolorbox.sty` from your project folder or local texmf tree.
% Fix: search your repo for `tcolorbox.sty` and delete/rename any local copy so TeX uses
% the TeX Live version. Also try "Clean" in LaTeX Workshop to remove stale .aux/.toc/.out.

\documentclass[12pt]{article}

% =========================
% 0) Metadata (EDIT ME)
% =========================
\newcommand{\CourseCode}{MATH 521}
\newcommand{\CourseTitle}{Analysis I}
\newcommand{\CourseSection}{LEC 002}
\newcommand{\Semester}{Spring 2026}
\newcommand{\Instructor}{Thierry Laurens}
\newcommand{\MeetingInfo}{MWF 1:20--2:10 \;|\; Van Vleck B119} % optional
\newcommand{\HomeworkNumber}{1} % <-- change each week
\newcommand{\YourName}{Alejandro De La Torre}
\newcommand{\DueDate}{\today} % <-- or set e.g. January 31, 2026

% =========================
% 1) Core math + proof tools
% =========================
\usepackage{amsmath, amssymb, amsthm}

% =========================
% 2) Lists and spacing
% =========================
\usepackage{enumitem}
\setlist[enumerate,1]{label=\textbf{(\alph*)}, leftmargin=2em, itemsep=0.25em}
\setlist[itemize]{leftmargin=1.5em, itemsep=0.25em}
\setlength{\parskip}{0.35em}
\setlength{\parindent}{0pt}

% =========================
% 3) Page geometry + header/footer
% =========================
\usepackage{geometry}
\geometry{margin=1in}

\usepackage{fancyhdr}
\pagestyle{fancy}
\fancyhf{}
\lhead{\CourseCode\ --- \CourseSection, \Semester \;(\MeetingInfo)}
\rhead{Homework \HomeworkNumber}
\cfoot{\thepage}
\renewcommand{\headrulewidth}{0.4pt}
% Fixes common fancyhdr warning about header height:
\setlength{\headheight}{15pt}
% Optional: reclaim a bit of vertical space after increasing headheight
\addtolength{\topmargin}{-2pt}

% =========================
% 4) Links (subtle)
% =========================
\usepackage[hidelinks]{hyperref}
\setcounter{tocdepth}{1}

% =========================
% 5) Quotes / figures
% =========================
\usepackage{csquotes}
\usepackage{graphicx}
\graphicspath{{images/}} % put figures in ./images/

% =========================
% 6) Simple scratch box (no tcolorbox)
% =========================
% A lightweight environment for brief planning / scratch work.
% This avoids any tcolorbox dependencies.
\newenvironment{scratch}{%
  \par\smallskip
  \noindent\textbf{Discussion \& Scratch Work}\begin{quote}\small
}{%
  \end{quote}\smallskip
}

% =========================
% 7) Lightweight theorem-like environments (optional)
% =========================
\newtheorem*{claim}{Claim}
\newtheorem*{lemma}{Lemma}
\newtheorem*{remark}{Remark}
\newtheorem{theorem}{Theorem}
\newtheorem{proposition}{Proposition}
\newtheorem{corollary}{Corollary}
\newtheorem{definition}{Definition}
\newtheorem{example}{Example}
\newtheorem{exercise}{Exercise}

% =========================
% 8) Problem header command (with TOC + PDF bookmarks)
% Usage: \problem[Optional short title]{<number>}
% =========================
\makeatletter
\newcommand{\problem}[2][]{%
  \def\prob@title{Problem #2\if\relax\detokenize{#1}\relax\else\ (\textit{#1})\fi}%
  \phantomsection%
  \pdfbookmark[1]{\prob@title}{prob#2}%
  \addcontentsline{toc}{section}{\prob@title}%
  \section*{\prob@title}%
  \vspace{-0.25em}\hrule\vspace{0.8em}%
}
\makeatother

% =========================
% 9) Title block
% =========================
\title{Homework \HomeworkNumber}
\author{\YourName \\
\CourseCode\ --- \CourseTitle \\
\CourseSection \\
Instructor: \Instructor}
\date{\DueDate}

\begin{document}
\maketitle

% ------------------ collaboration + sources ------------------
% \section*{Collaboration Statement}
% Write “I worked alone” or list collaborators / external resources.
% "I worked on this homework independently. No outside collaborators or resources were used."

\tableofcontents
\bigskip
\hrule
\bigskip

\newpage 

% ============================================================
% Problems (duplicate blocks as needed)
% ============================================================

\problem[]{1}
Using induction, prove that for all $n \in \mathbb{N}$ we have
\[
  3 + 11 + \cdots + (8n - 5) = 4n^2 - n.
\]

\textbf{Proof.}\\
For $n\in\mathbb{N}$, let $P(n)$ denote the statement
\[
  3 + 11 + \cdots + (8n-5) = 4n^2 - n.
\]
We argue by induction.

\emph{Base case.} For $n=1$, the left-hand side is $3$, while the right-hand side is $4\cdot 1^2-1=3$. Hence $P(1)$ holds.

\emph{Inductive step.} Assume $P(n)$ holds for some $n\in\mathbb{N}$, i.e.
\[
  3 + 11 + \cdots + (8n-5) = 4n^2 - n.
\]
We must show that $P(n+1)$ holds. Observe that the next term in the sum is
\[
  8(n+1)-5 = 8n+3.
\]
Adding this term to both sides of the inductive hypothesis gives
\[
  3 + 11 + \cdots + (8n-5) + (8n+3) = (4n^2-n) + (8n+3).
\]
The left-hand side is exactly the sum $3+11+\cdots + (8(n+1)-5)$, and the right-hand side simplifies to
\[
  (4n^2-n) + (8n+3) = 4n^2 + 7n + 3.
\]
On the other hand,
\[
  4(n+1)^2-(n+1) = 4(n^2+2n+1) - n - 1 = 4n^2 + 7n + 3.
\]
Therefore
\[
  3 + 11 + \cdots + (8(n+1)-5) = 4(n+1)^2-(n+1),
\]
so $P(n+1)$ holds. By the Principle of Mathematical Induction, $P(n)$ is true for all $n\in\mathbb{N}$.
\qed

\newpage

\problem[]{2}
Compute the following sum:
\[
  \frac{1}{1\cdot 4} + \frac{1}{4\cdot 7} + \cdots + \frac{1}{(3n-2)(3n+1)}.
\]
Prove that your answer is true for all $n \in \mathbb{N}$ using induction.

\begin{scratch}
We first compute the sum in closed form. For $k \in \mathbb{N}$,
\[
  \frac{1}{(3k-2)(3k+1)} = \frac{A}{3k-2} + \frac{B}{3k+1}.
\]
Solving
\[
  1 = A(3k+1) + B(3k-2)
\]
gives $A=\tfrac13$ and $B=-\tfrac13$. Hence
\[
  \frac{1}{(3k-2)(3k+1)}
  = \frac13\!\left(\frac{1}{3k-2} - \frac{1}{3k+1}\right).
\]

Therefore the partial sums telescope:
\[
  \sum_{k=1}^n \frac{1}{(3k-2)(3k+1)}
  = \frac13\left(1 - \frac{1}{3n+1}\right)
  = \frac{n}{3n+1}.
\]
We now prove this identity by induction.
\end{scratch}

\textbf{Proof.}\\
For $n \in \mathbb{N}$, let $P(n)$ denote the statement
\[
  \frac{1}{1\cdot 4} + \frac{1}{4\cdot 7} + \cdots + \frac{1}{(3n-2)(3n+1)}
  = \frac{n}{3n+1}.
\]
We will prove that $P(n)$ is true for all $n \in \mathbb{N}$ using induction.

\emph{Base case.}  
$P(1)$ is the statement
\[
  \frac{1}{1\cdot 4} = \frac{1}{3\cdot 1 + 1},
\]
which is true.

\emph{Inductive step.}  
Assume $P(n)$ is true for some $n \in \mathbb{N}$, that is,
\[
  \frac{1}{1\cdot 4} + \frac{1}{4\cdot 7} + \cdots + \frac{1}{(3n-2)(3n+1)}
  = \frac{n}{3n+1}.
\]
Adding the next term $\dfrac{1}{(3n+1)(3n+4)}$ to both sides, we obtain
\[
  \frac{1}{1\cdot 4} + \cdots + \frac{1}{(3n-2)(3n+1)} + \frac{1}{(3n+1)(3n+4)}
  = \frac{n}{3n+1} + \frac{1}{(3n+1)(3n+4)}.
\]
The left-hand side is exactly the sum in $P(n+1)$. On the right-hand side,
\[
  \frac{n}{3n+1} + \frac{1}{(3n+1)(3n+4)}
  = \frac{1}{3n+1}\left(n + \frac{1}{3n+4}\right)
  = \frac{1}{3n+1}\cdot\frac{n(3n+4)+1}{3n+4}.
\]
Since
\[
  n(3n+4)+1 = 3n^2+4n+1 = (3n+1)(n+1),
\]
we obtain
\[
  \frac{n}{3n+1} + \frac{1}{(3n+1)(3n+4)}
  = \frac{n+1}{3n+4}
  = \frac{n+1}{3(n+1)+1}.
\]
Thus $P(n+1)$ is true. By the Principle of Mathematical Induction, $P(n)$ is true for all $n \in \mathbb{N}$.
\qed

\newpage

\problem[]{3}
Let $f : A \to B$ be a function.
\begin{enumerate}
  \item Prove that $f$ is injective if and only if $f^{-1}(f(X)) = X$ for all $X \subseteq A$.
  \item Prove that $f$ is surjective if and only if $f\bigl(f^{-1}(Y)\bigr) = Y$ for all $Y \subseteq B$.
\end{enumerate}

\begin{scratch}
\end{scratch}

\textbf{Proof.}\\
\qed

\newpage

\problem[]{4}
Let $f : A \to B$ be a function. Prove that the following two statements are equivalent:
\begin{enumerate}
  \item The function $f$ is surjective.
  \item For any set $C$ and for any functions $g : B \to C$ and $h : B \to C$ such that $g \circ f = h \circ f$,
  we have $g = h$.
\end{enumerate}

\begin{scratch}
\end{scratch}

\textbf{Proof.}\\
\qed

\newpage

\problem[]{5}
Let $f : B \to C$ be a function. Prove that the following two statements are equivalent:
\begin{enumerate}
  \item The function $f$ is injective.
  \item For any set $A$ and for any functions $g : A \to B$ and $h : A \to B$ such that $f \circ g = f \circ h$,
  we have $g = h$.
\end{enumerate}

\begin{scratch}
\end{scratch}

\textbf{Proof.}\\
\qed

\newpage

\problem[]{6}
Let $f : A \to B$ be a function. Suppose that there exists a function $g : B \to A$
such that $(g \circ f)(a) = a$ for all $a \in A$ and $(f \circ g)(b) = b$ for all $b \in B$. Prove that $f$ is
bijective.

\begin{scratch}
\end{scratch}

\textbf{Proof.}\\
\qed

% ============================================================
% Appendix: Toolbox (Definitions & Theorems used in HW1)
% ============================================================
\newpage
\appendix
\section*{Appendix: Toolbox}
\addcontentsline{toc}{section}{Appendix: Toolbox}

\begin{definition}[Natural Numbers]
The set of natural numbers $\mathbb{N} = \{1,2,3,\dots\}$ is defined axiomatically by the Peano axioms:
\begin{enumerate}
  \item $1 \in \mathbb{N}$.
  \item If $n \in \mathbb{N}$, then $n+1 \in \mathbb{N}$.
  \item $1$ is not the successor of any natural number.
  \item If $n+1 = m+1$, then $n=m$.
  \item (\emph{Induction Axiom})  
  If $S \subseteq \mathbb{N}$ satisfies $1 \in S$ and $n \in S \Rightarrow n+1 \in S$, then $S=\mathbb{N}$.
\end{enumerate}
\end{definition}

\begin{theorem}[Principle of Mathematical Induction]
Let $P(n)$ be a statement defined for all $n \in \mathbb{N}$.  
If:
\begin{enumerate}
  \item $P(1)$ is true (base case), and
  \item for all $n \in \mathbb{N}$, $P(n) \Rightarrow P(n+1)$ (inductive step),
\end{enumerate}
then $P(n)$ is true for all $n \in \mathbb{N}$.
\end{theorem}

\begin{remark}
The Principle of Mathematical Induction follows from the Peano axioms.
Given a statement $P(n)$, define
\[
S = \{ n \in \mathbb{N} : P(n) \text{ is true} \}.
\]
If $1 \in S$ and $n \in S \Rightarrow n+1 \in S$, then $S=\mathbb{N}$ by the Induction Axiom.
\end{remark}

\begin{definition}[Function]
Let $A$ and $B$ be nonempty sets.  
A \emph{function} $f : A \to B$ assigns to each element $a \in A$ exactly one element $f(a) \in B$.
The set $A$ is called the \emph{domain} and the set $B$ the \emph{codomain}.
\end{definition}

\begin{definition}[Composition of Functions]
Let $f : A \to B$ and $g : B \to C$ be functions.  
The \emph{composition} of $g$ with $f$ is the function $g \circ f : A \to C$ defined by
\[
(g \circ f)(a) = g(f(a)) \quad \text{for all } a \in A.
\]
\end{definition}

\begin{definition}[Image and Preimage]
Let $f : A \to B$ be a function.
\begin{itemize}
  \item For $X \subseteq A$, the \emph{image} of $X$ under $f$ is
  \[
    f(X) = \{ f(x) : x \in X \}.
  \]
  \item For $Y \subseteq B$, the \emph{preimage} of $Y$ under $f$ is
  \[
    f^{-1}(Y) = \{ a \in A : f(a) \in Y \}.
  \]
\end{itemize}
\end{definition}

\begin{definition}[Injective, Surjective, Bijective]
Let $f : A \to B$ be a function.
\begin{itemize}
  \item $f$ is \emph{injective} if for all $a_1,a_2 \in A$,
  \[
    f(a_1)=f(a_2) \Rightarrow a_1=a_2.
  \]
  \item $f$ is \emph{surjective} if $f(A)=B$, i.e.\ for every $b \in B$ there exists $a \in A$ such that $f(a)=b$.
  \item $f$ is \emph{bijective} if it is both injective and surjective.
\end{itemize}
\end{definition}

\begin{remark}
For a general function $f : A \to B$, the notation $f^{-1}(Y)$ denotes a \emph{preimage} and does not require $f$ to be invertible.
Moreover, the properties of being injective, surjective, or bijective depend not only on the rule defining $f$,
but also on the choice of domain $A$ and codomain $B$.
\end{remark}

\end{document}